%%%%%%%%%%%%%%%%%%%%%%%%%%%%%%%%%%%%%%%%%%%%%%%%%%%%%%%%%%%%%%%%%
% GIÁO ÁN STEM: HÀM SỐ MŨ VÀ MÔ HÌNH TĂNG TRƯỞNG DỊCH BỆNH
% Môn: Toán 11 | Chương: Hàm số mũ và Logarit
% Phiên bản 2.0: Tiếp cận từ Toán học → Ứng dụng thực tiễn
%%%%%%%%%%%%%%%%%%%%%%%%%%%%%%%%%%%%%%%%%%%%%%%%%%%%%%%%%%%%%%%%%
\documentclass[12pt,a4paper]{article}

% --- GÓI CƠ BẢN ---
\usepackage[utf8]{inputenc}
\usepackage[T1]{fontenc}
\usepackage[vietnamese]{babel}
\usepackage{graphicx}
\usepackage{amsmath,amssymb}
\usepackage{array}
\usepackage{booktabs}
\usepackage{multirow}
\usepackage{tabularx}
\usepackage{colortbl}
\usepackage{wrapfig}
\usepackage{ragged2e}
\usepackage{setspace}
\usepackage{tikz}
\usetikzlibrary{shapes,arrows,positioning,patterns,calc}
\usepackage{fontawesome5}
\usepackage{verbatim}

% --- FONT ---
\usepackage{newtxtext, newtxmath}
\usepackage[scaled=0.92]{helvet}
\usepackage{courier}
\usepackage{microtype}

% --- MÀU SẮC ---
\usepackage[svgnames]{xcolor}
\definecolor{primary}{RGB}{0,102,153}
\definecolor{secondary}{RGB}{51,153,102}
\definecolor{accent}{RGB}{255,102,0}
\definecolor{mathblue}{RGB}{0,51,102}
\definecolor{success}{RGB}{34,139,34}
\definecolor{TableHeader}{gray}{0.9}

\usepackage[left=2cm,right=2cm,top=2.5cm,bottom=2cm,headheight=15.6pt]{geometry}
\usepackage{fancyhdr}
\pagestyle{fancy}
\fancyhf{}
\renewcommand{\headrule}{\color{primary}\hrule width\textwidth height 1.5pt}
\fancyhead[L]{\color{primary}\footnotesize\sffamily\textbf{SỞ GD\&ĐT TP.HCM - HỘI THI STEM/STEAM CẤP THÀNH PHỐ}}
\fancyhead[R]{\color{primary}\footnotesize\sffamily\textbf{NĂM HỌC 2025-2026}}
\fancyfoot[C]{\sffamily\thepage}
\renewcommand{\footrule}{\color{primary}\hrule width\textwidth height 0.8pt}

% --- GÓI HỖ TRỢ ---
\usepackage{enumitem}
\usepackage{tcolorbox}
\tcbuselibrary{skins, breakable}
\usepackage[colorlinks=true,linkcolor=primary,urlcolor=accent,citecolor=success,
    pdftitle={Giao an STEM: Ham so Mu va Mo hinh Dich benh},pdfauthor={Nguyen Van Sang}]{hyperref}

% --- TIÊU ĐỀ ---
\usepackage{titlesec}
\titleformat{\subsection}{\normalfont\bfseries\color{accent}\sffamily}
{\llap{\color{accent}\rule{1.2em}{1.2em}\hspace{0.5em}}}{0em}{}
\titleformat{\subsubsection}{\normalfont\bfseries\itshape\color{secondary}\sffamily}
{\llap{\color{secondary}\rule{1em}{1em}\hspace{0.5em}}}{0em}{}

% --- MÔI TRƯỜNG ---
\def\phanI{TỔNG QUAN}
\def\phanII{CHUẨN BỊ CỦA GIÁO VIÊN}
\def\phanIII{KẾ HOẠCH DẠY HỌC ĐỀ XUẤT}
\def\phanIV{HƯỚNG DẪN HỌC SINH}
\def\phanV{CÁC PHỤ LỤC}

\newtcolorbox{competitionsection}[1]{
    enhanced, breakable, colback=white, colframe=primary, boxrule=0.8pt,
    fonttitle=\bfseries\sffamily\Huge, coltitle=white, colbacktitle=primary,
    attach boxed title to top left={xshift=1cm, yshift=-1mm-\tcboxedtitleheight/2},
    boxed title style={sharp corners, boxrule=0pt},
    title={PHẦN #1: \csname phan#1\endcsname},
    before upper={\addcontentsline{toc}{section}{Phần #1: \csname phan#1\endcsname}}
}

\newcolumntype{C}{>{\centering\arraybackslash}X}
\newcolumntype{L}{>{\raggedright\arraybackslash}X}
\newcommand{\header}[1]{\cellcolor{TableHeader}\sffamily\bfseries#1}

\newtcolorbox{mathbox}{enhanced, breakable, colback=mathblue!5, colframe=mathblue!80, boxrule=1.5pt, arc=2mm,
    title=\sffamily\bfseries\color{white}KIẾN THỨC TOÁN HỌC, fonttitle=\bfseries\sffamily, colbacktitle=mathblue}
\newtcolorbox{realworld}{enhanced, breakable, colback=accent!5, colframe=accent!80, boxrule=1pt, arc=2mm,
    title=\sffamily\bfseries\color{primary}ỨNG DỤNG THỰC TIỄN, fonttitle=\bfseries\sffamily}
\newtcolorbox{task}[1]{enhanced, breakable, colback=secondary!10, colframe=secondary!80, boxrule=1pt, arc=2mm,
    title=\sffamily\bfseries\color{white}NHIỆM VỤ: #1, fonttitle=\bfseries\sffamily, colbacktitle=secondary}

\setlist[itemize,1]{label=\textbullet, leftmargin=*, nosep}
\setlist[itemize,2]{label=--, leftmargin=*, nosep}
\setlist[description]{font=\bfseries\sffamily\color{primary}, nosep}

\begin{document}

% --- TRANG BÌA ---
\begin{titlepage}
\begin{tikzpicture}[remember picture, overlay]
\node[rectangle, fill=primary, anchor=north, minimum width=\paperwidth, minimum height=2cm] (header) at (current page.north){};
\node[anchor=south, yshift=2mm, text=white] at (header.south) {\Huge\sffamily\bfseries GIÁO ÁN DỰ THI STEM/STEAM CẤP THÀNH PHỐ};
\end{tikzpicture}

\vspace{2cm}
\begin{center}
{\LARGE\bfseries\color{primary}\sffamily SỞ GIÁO DỤC VÀ ĐÀO TẠO THÀNH PHỐ HỒ CHÍ MINH}\\
\vspace{0.5cm}
{\Large\bfseries\sffamily HỘI THI THIẾT KẾ CHỦ ĐỀ DẠY HỌC STEM/STEAM}\\
\vspace{0.3cm}
{\large\sffamily Năm học 2025 - 2026}

\vspace{2cm}

\begin{tcolorbox}[width=0.9\textwidth, colback=gray!5, colframe=primary, arc=2mm, boxrule=1pt]
    \centering
    {\LARGE\bfseries\sffamily CHỦ ĐỀ:} \\[0.3cm]
    {\Huge\bfseries\color{mathblue}\sffamily HÀM SỐ MŨ}\\
    \vspace{0.2cm}
    {\Large\sffamily VÀ MÔ HÌNH TĂNG TRƯỞNG DỊCH BỆNH}\\
    \vspace{0.3cm}
    \rule{0.6\textwidth}{0.5pt}\\
    \vspace{0.2cm}
    {\large\itshape\sffamily Từ Toán học đến Ứng dụng Y tế Công cộng}
\end{tcolorbox}

\vspace{1.5cm}
{\Large\bfseries\sffamily MÔN: TOÁN - KHỐI: 11}\\
\vspace{0.3cm}
{\large\sffamily Chương: Hàm số mũ và Logarit}

\vfill
\large
\begin{tabular}{rl}
\sffamily\bfseries Giáo viên: & Nguyễn Văn Sang\\
\sffamily\bfseries Đơn vị: & Trường THPT Nguyễn Hữu Cảnh\\
\sffamily\bfseries Số điện thoại: & 0389\,821\,115\\
\end{tabular}

\vspace{1cm}
{\large\sffamily Thành phố Hồ Chí Minh, tháng 11 năm 2025}
\end{center}
\end{titlepage}

\tableofcontents
\thispagestyle{empty}

%%%%% PHẦN I: TỔNG QUAN %%%%%
\begin{competitionsection}{I}
\vspace{0.5cm}
\subsection*{1.1. Thông tin cá nhân}
\begin{itemize}[leftmargin=*,noitemsep]
\item \textbf{Họ và tên:} Nguyễn Văn Sang
\item \textbf{Chuyên môn:} Toán học
\item \textbf{Thâm niên:} 15 năm giảng dạy THPT
\item \textbf{Đơn vị:} Trường THPT Nguyễn Hữu Cảnh, TP.HCM
\end{itemize}

\subsection*{1.2. Giới thiệu chủ đề}
\begin{itemize}[leftmargin=*,noitemsep]
\item \textbf{Tên chủ đề:} \textcolor{mathblue}{Hàm số Mũ và Mô hình Tăng trưởng Dịch bệnh}
\item \textbf{Xuất phát điểm:} Từ \textbf{kiến thức Toán học} (Hàm số mũ) đến \textbf{ứng dụng thực tiễn}
\item \textbf{Bối cảnh Toán học:}
    \begin{itemize}
    \item Chương trình Toán 11 CTGDPT 2018: \textbf{Hàm số mũ và Logarit}
    \item Hàm số mũ $y = a^x$ với $a > 0, a \neq 1$ là nền tảng cho nhiều mô hình tăng trưởng
    \item Học sinh thường hỏi: \textit{``Học hàm số mũ để làm gì?''}
    \end{itemize}
\item \textbf{Điểm mới và sáng tạo:}
    \begin{itemize}
    \item \textbf{Toán học là trục chính:} Bắt đầu từ công thức $y = y_0 \cdot e^{rt}$, sau đó ứng dụng
    \item \textbf{Đa dạng ứng dụng:} Lãi kép, tăng trưởng dân số, phân rã phóng xạ, \textbf{và dịch bệnh}
    \item \textbf{Công nghệ hỗ trợ:} Ứng dụng web mô phỏng trực quan
    \item \textbf{Giải đáp câu hỏi:} Cho học sinh thấy Toán học có ích trong đời sống
    \end{itemize}
\end{itemize}

\subsection*{1.3. Căn cứ thiết kế}
\begin{minipage}{0.48\textwidth}
\textbf{Căn cứ pháp lý:}
\begin{itemize}[leftmargin=*,noitemsep]
\item Chương trình GDPT 2018 môn Toán lớp 11
\item Nội dung: Hàm số mũ và logarit
\item Công văn 5555/BGDĐT-GDTrH về dạy học STEM
\end{itemize}
\end{minipage}
\hfill
\begin{minipage}{0.48\textwidth}
\textbf{Yêu cầu cần đạt (SGK Toán 11):}
\begin{itemize}[leftmargin=*,noitemsep]
\item Nhận biết hàm số mũ và tính chất
\item Vẽ đồ thị hàm số mũ
\item Vận dụng vào bài toán thực tế
\end{itemize}
\end{minipage}

\subsection*{1.4. Mục tiêu giáo dục}
\begin{center}
\begin{tabularx}{\linewidth}{@{} L L L @{}}
\toprule
\header{Kiến thức Toán học} & \header{Kỹ năng STEM} & \header{Thái độ} \\
\midrule
\parbox[t]{\dimexpr1.\linewidth-2\tabcolsep}{
    \begin{itemize}[leftmargin=*,noitemsep,topsep=1ex]
        \item Hiểu hàm số mũ $y = a^x$
        \item Viết công thức tăng trưởng $y = y_0 \cdot e^{rt}$
        \item Tính toán với logarit tự nhiên
    \end{itemize}
} &
\parbox[t]{\dimexpr1.\linewidth-2\tabcolsep}{
    \begin{itemize}[leftmargin=*,noitemsep,topsep=1ex]
        \item Mô hình hóa bài toán thực tế
        \item Sử dụng công nghệ mô phỏng
        \item Đọc hiểu đồ thị, phân tích dữ liệu
    \end{itemize}
} &
\parbox[t]{\dimexpr1.\linewidth-2\tabcolsep}{
    \begin{itemize}[leftmargin=*,noitemsep,topsep=1ex]
        \item Thấy được ý nghĩa của Toán học
        \item Hứng thú với ứng dụng thực tế
        \item Tư duy phản biện
    \end{itemize}
} \\
\bottomrule
\end{tabularx}
\end{center}

\subsection*{1.5. Cấu trúc chương trình: Từ Toán đến Thực tiễn}
\begin{center}
\begin{tikzpicture}[node distance=1.2cm, every node/.style={font=\sffamily}]
\node (math1) [rectangle, rounded corners, fill=mathblue!20, text width=5cm, align=center, minimum height=1cm] {\textbf{TOÁN HỌC}\\Hàm số mũ $y = a^x$};
\node (math2) [rectangle, rounded corners, fill=mathblue!30, text width=5cm, align=center, below of=math1, yshift=-0.8cm, minimum height=1cm] {\textbf{CÔNG THỨC}\\$y = y_0 \cdot e^{rt}$};
\node (app1) [rectangle, rounded corners, fill=accent!20, text width=5cm, align=center, below of=math2, yshift=-0.8cm, minimum height=1cm] {\textbf{ỨNG DỤNG 1}\\Lãi kép ngân hàng};
\node (app2) [rectangle, rounded corners, fill=accent!30, text width=5cm, align=center, below of=app1, yshift=-0.8cm, minimum height=1cm] {\textbf{ỨNG DỤNG 2}\\Mô hình dịch bệnh (SIR)};
\draw[->, thick, color=mathblue] (math1) -- (math2);
\draw[->, thick, color=accent] (math2) -- (app1);
\draw[->, thick, color=accent] (app1) -- (app2);
\node[right of=math1, xshift=3cm, text width=4cm] {\footnotesize Định nghĩa, tính chất, đồ thị};
\node[right of=math2, xshift=3cm, text width=4cm] {\footnotesize Mô hình tăng trưởng liên tục};
\node[right of=app1, xshift=3cm, text width=4cm] {\footnotesize $A = P \cdot e^{rt}$};
\node[right of=app2, xshift=3cm, text width=4cm] {\footnotesize $I(t) = I_0 \cdot e^{(\beta-\gamma)t}$};
\end{tikzpicture}
\end{center}
\end{competitionsection}
\newpage

%%%%% PHẦN II: CHUẨN BỊ %%%%%
\begin{competitionsection}{II}
\vspace{0.5cm}
\subsection*{2.1. Thiết bị và công nghệ}
\begin{itemize}[leftmargin=*,noitemsep]
\item \textbf{Phần cứng:} Máy tính/điện thoại có trình duyệt web
\item \textbf{Phần mềm:}
    \begin{itemize}
    \item \textbf{Học liệu số:} Ứng dụng web \texttt{EpidemicSim.html} hoặc \href{https://stem-1h8.pages.dev/EpidemicSim.html}{stem-1h8.pages.dev/EpidemicSim.html}
    \item Máy tính cầm tay Casio fx-580VN X (tính $e^x$, $\ln x$)
    \end{itemize}
\item \textbf{Tài liệu:}
    \begin{itemize}
    \item Phiếu học tập số 1: Khám phá Hàm số mũ
    \item Phiếu học tập số 2: Ứng dụng vào Mô hình dịch bệnh
    \item Rubric đánh giá sản phẩm
    \end{itemize}
\end{itemize}

\subsection*{2.2. Kịch bản giảng dạy (2 tiết = 90 phút)}
\begin{center}
\begin{tabularx}{\linewidth}{@{} L L L L @{}}
\toprule
\header{Thời gian} & \header{Nội dung} & \header{Phương pháp} & \header{Công cụ} \\
\midrule
10' & Khởi động: Câu hỏi ``Học hàm mũ để làm gì?'' & Đàm thoại & Slide câu hỏi \\
\rowcolor{mathblue!10}
20' & Ôn tập Hàm số mũ: $y = a^x$, tính chất, đồ thị & Giảng dạy trực tiếp & Bảng, PHT1 \\
\rowcolor{mathblue!10}
10' & Công thức tăng trưởng: $y = y_0 \cdot e^{rt}$ & Khám phá có hướng dẫn & PHT1 \\
\rowcolor{accent!10}
10' & Ứng dụng 1: Lãi kép ngân hàng & Bài toán thực tế & Máy tính \\
\rowcolor{accent!10}
25' & Ứng dụng 2: Mô hình dịch bệnh & Học qua mô phỏng & EpidemicSim \\
10' & Trình bày và phản biện & Thuyết trình nhóm & Slides \\
5' & Tổng kết: Toán học trong đời sống & Suy ngẫm & -- \\
\bottomrule
\end{tabularx}
\end{center}

\subsection*{2.3. Điểm nhấn: Tiếp cận từ Toán học}
\begin{mathbox}
\textbf{Hàm số mũ cơ sở $e$:}
$$y = e^x \quad \text{với} \quad e = \lim_{n \to \infty}\left(1 + \frac{1}{n}\right)^n \approx 2.71828...$$

\textbf{Mô hình tăng trưởng liên tục:}
$$y(t) = y_0 \cdot e^{rt}$$
\begin{itemize}[leftmargin=*,noitemsep]
\item $y_0$ = Giá trị ban đầu (tại $t=0$)
\item $r$ = Tốc độ tăng trưởng (nếu $r > 0$: tăng, $r < 0$: giảm)
\item $t$ = Thời gian
\end{itemize}

\textbf{Ý nghĩa logarit:} Tìm thời gian $t$ khi biết $y$:
$$t = \frac{\ln(y/y_0)}{r}$$
\end{mathbox}

\begin{realworld}
\textbf{Các ứng dụng của công thức $y = y_0 \cdot e^{rt}$:}
\begin{enumerate}[leftmargin=*,noitemsep]
\item \textbf{Lãi kép ngân hàng:} $A = P \cdot e^{rt}$ (P: vốn gốc, r: lãi suất, t: thời gian)
\item \textbf{Tăng trưởng dân số:} $N = N_0 \cdot e^{rt}$ (r: tỷ lệ sinh - tỷ lệ tử)
\item \textbf{Phân rã phóng xạ:} $N = N_0 \cdot e^{-\lambda t}$ (r = $-\lambda$ < 0: giảm)
\item \textbf{Lây lan dịch bệnh:} $I = I_0 \cdot e^{(\beta - \gamma)t}$ (giai đoạn đầu)
\end{enumerate}
\end{realworld}
\end{competitionsection}

%%%%% PHẦN III: KẾ HOẠCH DẠY HỌC %%%%%
\begin{competitionsection}{III}
\vspace{0.5cm}
\subsection*{3.1. Kịch bản chi tiết}

\subsubsection*{Hoạt động 1: Khởi động - Đặt vấn đề (10 phút)}
\begin{itemize}[leftmargin=*,noitemsep]
\item \textbf{Câu hỏi mở đầu:} \textit{``Học hàm số mũ để làm gì trong cuộc sống?''}
\item \textbf{Học sinh thảo luận:} Thu thập ý kiến (lãi suất, dân số, virus...)
\item \textbf{Giáo viên dẫn dắt:} \textit{``Hôm nay chúng ta sẽ khám phá sức mạnh của hàm số mũ qua nhiều ứng dụng thực tế, đặc biệt là mô hình lây lan dịch bệnh!''}
\end{itemize}

\subsubsection*{Hoạt động 2: Ôn tập Hàm số Mũ (20 phút)}
\begin{task}{Cá nhân + Nhóm}
\textbf{Phiếu học tập số 1:} Ôn tập kiến thức Hàm số mũ
\begin{itemize}[leftmargin=*,noitemsep]
\item Định nghĩa hàm số mũ $y = a^x$ ($a > 0$, $a \neq 1$)
\item Tính chất: $a^0 = 1$, $a^{x+y} = a^x \cdot a^y$, ...
\item Đồ thị: qua điểm $(0, 1)$, tăng nếu $a > 1$, giảm nếu $0 < a < 1$
\item \textbf{Số $e$:} $e \approx 2.71828$, hàm $y = e^x$ là ``hàm mũ tự nhiên''
\end{itemize}
\end{task}

\subsubsection*{Hoạt động 3: Công thức Tăng trưởng (10 phút)}
\begin{mathbox}
\textbf{Bài toán:} Một đại lượng $y$ tăng trưởng liên tục với tốc độ $r$ (r\% mỗi đơn vị thời gian).

\textbf{Công thức:}
$$y(t) = y_0 \cdot e^{rt}$$

\textbf{Ví dụ 1:} Gửi ngân hàng 100 triệu, lãi suất 5\%/năm, lãi kép liên tục. Sau 10 năm:
$$A = 100 \cdot e^{0.05 \times 10} = 100 \cdot e^{0.5} \approx 164.87 \text{ triệu}$$

\textbf{Ví dụ 2:} Dân số một thành phố ban đầu 1 triệu, tăng 2\%/năm. Sau 20 năm:
$$N = 1 \cdot e^{0.02 \times 20} = e^{0.4} \approx 1.492 \text{ triệu}$$
\end{mathbox}

\subsubsection*{Hoạt động 4: Ứng dụng vào Mô hình Dịch bệnh (25 phút)}
\begin{realworld}
\textbf{Tình huống:} Một thành phố có 10,000 dân. Xuất hiện 10 ca nhiễm bệnh X ban đầu.

\textbf{Câu hỏi:} Dùng công thức tăng trưởng mũ để dự đoán số ca nhiễm theo thời gian?

\textbf{Thông số:}
\begin{itemize}[leftmargin=*,noitemsep]
\item Tốc độ lây lan: $r = \beta - \gamma = 0.3 - 0.1 = 0.2$ (20\%/ngày)
\item Công thức: $I(t) = 10 \cdot e^{0.2t}$
\end{itemize}
\end{realworld}

\begin{task}{Nhóm - Sử dụng EpidemicSim}
\textbf{Bước 1:} Mở ứng dụng web \texttt{EpidemicSim.html}

\textbf{Bước 2:} Nhập tham số ban đầu:
\begin{itemize}[leftmargin=*,noitemsep]
\item Dân số N = 10,000
\item Số ca ban đầu $I_0$ = 10
\item $\beta$ = 0.3, $\gamma$ = 0.1
\end{itemize}

\textbf{Bước 3:} Quan sát đồ thị và trả lời:
\begin{enumerate}[leftmargin=*,noitemsep]
\item Giai đoạn đầu có dạng hàm mũ không? (so sánh với $I = 10 \cdot e^{0.2t}$)
\item Tại sao sau một thời gian đồ thị không còn tăng theo hàm mũ?
\item Đỉnh dịch (peak) xảy ra khi nào? Bao nhiêu ca?
\end{enumerate}

\textbf{Bước 4:} Thử nghiệm: Thay đổi $\beta$ (giãn cách xã hội) hoặc thêm tiêm chủng.
\end{task}

\subsubsection*{Hoạt động 5: Trình bày và Phản biện (10 phút)}
\begin{itemize}[leftmargin=*,noitemsep]
\item 2 nhóm trình bày kết quả mô phỏng và nhận xét
\item Các nhóm khác đặt câu hỏi phản biện
\item Giáo viên tổng kết: Liên hệ Toán học với thực tiễn
\end{itemize}

\subsubsection*{Hoạt động 6: Tổng kết (5 phút)}
\begin{itemize}[leftmargin=*,noitemsep]
\item \textbf{Kiến thức Toán:} Hàm số mũ $y = a^x$, công thức $y = y_0 \cdot e^{rt}$
\item \textbf{Ứng dụng:} Lãi kép, dân số, dịch bệnh
\item \textbf{Kỹ năng STEM:} Mô hình hóa, sử dụng công nghệ, phân tích dữ liệu
\item \textbf{Bài tập về nhà:}
    \begin{enumerate}[leftmargin=*,noitemsep]
    \item Tính số tiền sau 5 năm nếu gửi 50 triệu với lãi suất 6\%/năm (lãi kép liên tục)
    \item Một chất phóng xạ có chu kỳ bán rã 10 ngày. Sau bao lâu còn 25\% khối lượng ban đầu?
    \end{enumerate}
\end{itemize}
\end{competitionsection}
\newpage

%%%%% PHẦN IV: HƯỚNG DẪN HỌC SINH %%%%%
\begin{competitionsection}{IV}
\vspace{0.5cm}
\subsection*{4.1. Tóm tắt kiến thức}

\begin{mathbox}
\textbf{1. Hàm số mũ:}
$$y = a^x \quad (a > 0, a \neq 1)$$

\textbf{2. Hàm số mũ cơ sở $e$:}
$$y = e^x \quad \text{với} \quad e \approx 2.71828$$

\textbf{3. Mô hình tăng trưởng liên tục:}
$$y(t) = y_0 \cdot e^{rt}$$

\textbf{4. Công thức nghịch đảo (dùng logarit):}
$$t = \frac{\ln(y/y_0)}{r} = \frac{\ln y - \ln y_0}{r}$$
\end{mathbox}

\subsection*{4.2. Ví dụ ứng dụng}

\textbf{Ví dụ 1: Lãi kép}\\
Gửi 100 triệu với lãi suất 5\%/năm, lãi kép liên tục. Sau $t$ năm:
$$A = 100 \cdot e^{0.05t}$$
Sau 10 năm: $A = 100 \cdot e^{0.5} \approx 164.87$ triệu.

\textbf{Ví dụ 2: Dịch bệnh (giai đoạn đầu)}\\
Ban đầu có 10 ca nhiễm, tốc độ lây lan $r = 0.2$/ngày:
$$I(t) = 10 \cdot e^{0.2t}$$
Sau 10 ngày: $I = 10 \cdot e^{2} \approx 73.9$ ca.\\
Sau 20 ngày: $I = 10 \cdot e^{4} \approx 545.9$ ca.

\textbf{Lưu ý:} Công thức mũ chỉ đúng ở giai đoạn đầu. Khi số người nhạy cảm giảm, tốc độ lây lan chậm lại (mô hình SIR chính xác hơn).

\subsection*{4.3. Hướng dẫn sử dụng EpidemicSim}
\begin{itemize}[leftmargin=*,noitemsep]
\item \textbf{Bước 1:} Truy cập \href{https://stem-1h8.pages.dev/EpidemicSim.html}{stem-1h8.pages.dev/EpidemicSim.html}
\item \textbf{Bước 2:} Tab ``Mô phỏng'' - Nhập tham số và chạy mô phỏng
\item \textbf{Bước 3:} So sánh đường cong I(t) với công thức $I_0 \cdot e^{rt}$ (giai đoạn đầu)
\item \textbf{Bước 4:} Tab ``So sánh'' - Thử các kịch bản khác nhau
\item \textbf{Bước 5:} Tab ``Học tập'' - Đọc thêm về mô hình SIR
\end{itemize}
\end{competitionsection}
\newpage

%%%%% PHẦN V: PHỤ LỤC %%%%%
\begin{competitionsection}{V}
\vspace{0.5cm}
\subsection*{Phụ lục 1: Bảng giá trị $e^x$}
\begin{center}
\begin{tabular}{|c|c|c|c|c|c|c|c|}
\hline
$x$ & 0 & 0.5 & 1 & 1.5 & 2 & 3 & 4 \\
\hline
$e^x$ & 1 & 1.649 & 2.718 & 4.482 & 7.389 & 20.09 & 54.60 \\
\hline
\end{tabular}
\end{center}

\subsection*{Phụ lục 2: So sánh các mô hình tăng trưởng}
\begin{center}
\begin{tabularx}{0.95\linewidth}{@{} L C C @{}}
\toprule
\header{Loại tăng trưởng} & \header{Công thức} & \header{Đặc điểm} \\
\midrule
Tuyến tính & $y = y_0 + rt$ & Tăng đều, chậm \\
Lũy thừa & $y = y_0 \cdot t^n$ & Tăng nhanh dần \\
Mũ & $y = y_0 \cdot e^{rt}$ & Tăng rất nhanh (bùng nổ) \\
Logistic & $y = \dfrac{K}{1 + Ae^{-rt}}$ & Tăng rồi bão hòa \\
\bottomrule
\end{tabularx}
\end{center}

\subsection*{Phụ lục 3: Liên hệ với Mô hình SIR}
\begin{itemize}[leftmargin=*,noitemsep]
\item \textbf{Mô hình SIR đầy đủ:} Gồm 3 nhóm S (Susceptible), I (Infected), R (Recovered)
\item \textbf{Giai đoạn đầu:} Khi $S \approx N$ (gần như tất cả đều nhạy cảm), ta có:
$$\frac{dI}{dt} \approx (\beta - \gamma) I \quad \Rightarrow \quad I(t) = I_0 \cdot e^{(\beta - \gamma)t}$$
\item Đây chính là \textbf{công thức tăng trưởng mũ} với $r = \beta - \gamma$
\item \textbf{Sau đó:} Số người nhạy cảm giảm, tốc độ lây lan chậm lại, đồ thị uốn cong và đạt đỉnh
\item \textbf{Hệ số $R_0 = \beta/\gamma$:} Nếu $R_0 > 1$ thì dịch bùng phát, $R_0 < 1$ thì dịch tự tắt
\end{itemize}

\subsection*{Phụ lục 4: Rubric đánh giá}
\begin{tabularx}{\textwidth}{|>{\raggedright\arraybackslash}p{2.5cm}|X|X|X|c|}
\hline
\rowcolor{TableHeader}
\sffamily\bfseries Tiêu chí & \sffamily\bfseries Xuất sắc (4đ) & \sffamily\bfseries Khá (3đ) & \sffamily\bfseries Cần cải thiện (1-2đ) & \sffamily\bfseries Điểm \\
\hline
\textbf{Kiến thức Toán} & Nắm vững hàm mũ, tính đúng $e^{rt}$ & Hiểu cơ bản, tính đúng & Còn nhầm lẫn công thức & \underline{\hspace{1cm}} \\
\hline
\textbf{Mô hình hóa} & Viết đúng công thức ứng dụng & Viết được với hỗ trợ & Chưa viết được & \underline{\hspace{1cm}} \\
\hline
\textbf{Sử dụng công nghệ} & Thành thạo mô phỏng, phân tích & Sử dụng cơ bản & Cần hướng dẫn & \underline{\hspace{1cm}} \\
\hline
\textbf{Trình bày} & Mạch lạc, có minh họa & Rõ ràng & Rời rạc & \underline{\hspace{1cm}} \\
\hline
\end{tabularx}
\vspace{0.5cm}
\noindent\textbf{TỔNG ĐIỂM:} \underline{\hspace{2cm}} / 16 điểm
\end{competitionsection}

\end{document}
